%# -*- coding: utf-8-unix -*-

\chapter*{前言}
\addcontentsline{toc}{chapter}{前言}

\begin{quote}
\centering\itshape
Fermat, Euler, Lagrange, Legendre... introitum ad penetralia huius divinae scientiae aperuerunt, quantisque divitiis abundent patefecerunt. \\
\hfill --- Gauss, \textit{Disquisitiones Arithmeticae}
\end{quote}

\vspace{1em}

超越数的研究源于多种不同的背景，例如古希腊关于化圆为方的难题、Liouville 和 Cantor 的初步研究、Hermite 对指数函数的研究，以及 Hilbert 著名的 23 个问题中的第七个问题. 如今，它已发展成为一个蓬勃发展且广泛的理论，丰富了数学的各个分支. 因此，现在正是撰写一部系统性著作的契机. 本书的目的旨在对该领域近期取得的重大发现提供全面的阐述. 具体而言，书中介绍了与代数数对数线性型相关的最新理论、Schmidt 对 Thue-Siegel-Roth 定理的推广、Shidlovsky 关于 Siegel 的 $E$-函数的工作，以及 Sprindzuk 对 Mahler 猜想的证明. 本书在阐述过程中也讨论了该学科的经典内容. 特别地，为了便于读者建立清晰的历史脉络，开头部分对世纪之交（约 20 世纪初）该理论的发展状况进行了概述. 可以理解的是，该学科中的证明往往冗长而复杂，因此有必要仅选择最基本的结论进行详细讨论. 此外，总体而言，重点放在了那些引出迄今为止已知最强命题，或产生最广泛应用的方法上. 尽管如此，仍希望书中已包含了对相关文献的充分引用.

尽管历史悠久，但显而易见的是，超越数理论依然充满着年轻的活力. 许多课题无疑将受益于更深一步的研究，而且一些著名的长期悬而未决的问题仍然未解. 例如，只需提及关于 $e$ 与 $\pi$ 的代数独立性，以及 Euler 常数 $\gamma$ 的超越性这些著名的猜想，其中任何一个问题的解决都将代表一个重大进展. 如果本书能在推动未来的研究进展中起到微薄的作用，作者便深感欣慰.

本书内容源于在 Cambridge、美国以及其他地方举办的多次讲座，同时也构成了 Adams 奖论文的主要内容.

我非常感谢 D. W. Masser 博士在校对证明方面给予的友好帮助，同时也感谢 Cambridge University Press 在印刷方面所付出的心血.

\vspace{1em}

\noindent
\textit{Cambridge, 1974} \hfill \textbf{A. B.}
